% Response-aware allocation section for tech report

This section documents an investigation into whether the coefficient-SSE
objective could be replaced by a frequency-response-aware objective, and
whether the resulting allocation could be solved exactly under exact
quota constraints.

\subsection{Motivation and Mathematical Structure}

The coefficient-SSE objective $\sum_i e_i^2$ is (by Parseval's theorem)
proportional to the unweighted integrated squared frequency-response
error over the full band. To focus on stopband compliance, the natural
objective is the stopband-weighted integrated squared complex-response
error $\mathbf{e}^\top R \mathbf{e}$, where
\begin{equation}
R_{ij} = \frac{1}{\pi} \int_{\Omega_\text{stop}} \cos(\omega(i-j))\,d\omega
\end{equation}
and $\Omega_\text{stop}$ is the stopband interval.

Three key properties were established:
\begin{enumerate}
\item $R_{ii}$ is the same for all tap positions $i$: the diagonal
weight equals the stopband bandwidth fraction and does not vary with
tap index. Diagonal approximation therefore produces a uniform
scaling of the SSE objective and does not change the allocation.
\item $R$ has nonzero off-diagonal elements (magnitudes $6\times10^{-6}$
to $0.30$ for the 48-tap case), so the objective $\mathbf{e}^\top R
\mathbf{e}$ is not separable in the tap errors $e_i$. The existing
quota DP, which assumes a separable objective, cannot directly solve
the weighted formulation.
\item Fang et al.~\cite{fang2025mixedprecision} also studied FIR
allocation with a frequency-domain objective (minimax), so the
objective type alone does not establish novelty.
\end{enumerate}

\subsection{Small-Scale Exhaustive Screening}

To assess whether handling the cross-terms changes the optimal
allocation, exhaustive search was run on five 8-tap filters under two
quotas each (10 cases). Each case enumerated all feasible assignments
(420--560 per case), found the SSE-minimizing allocation and the
$\mathbf{e}^\top R \mathbf{e}$-minimizing allocation, and compared them.

In 6 of 10 cases, the two optima differed and the weighted-optimal
allocation had a lower stopband peak. However, the absolute stopband
peaks for these 8-tap filters are poor ($-5$ to $-20$\,dB), and the
improvements are small ($0.2$--$1.2$\,dB). The 10 cases arise from 5
filter designs $\times$ 2 quotas; independence is limited.

\subsection{Forty-Eight-Tap MIQP Attempt}

To test whether the weighted objective can be solved exactly at realistic
filter scales, the objective was linearized using binary product
variables. For each pair $(i,k),(j,l)$ with $i<j$, a product variable
$y = x_{i,k} x_{j,l}$ was introduced with standard McCormick
inequalities. The resulting MILP had:
\begin{itemize}
\item 144 width-selection binary variables
\item 10,152 product binary variables
\item 51 equality constraints
\item 30,456 McCormick inequality constraints
\end{itemize}

MATLAB \texttt{intlinprog} was warm-started with the SSE-DP solution.
After approximately 93 seconds, the root-node LP lower bound was
approximately $-5.637 \times 10^{-4}$ against an incumbent of
$7.4603 \times 10^{-6}$ (gap~$\approx 7{,}656\%$), with no
branch-and-bound nodes explored. The run was manually interrupted after
approximately 93 seconds under the pre-declared slow-progress stopping
condition. The root bound changed negligibly from approximately
$-5.6371\times10^{-4}$. Ctrl+C terminated MATLAB before it returned an
exit flag or final allocation. The negative LP bound reflects weak LP
relaxation from the McCormick linearization; it is not a property of the
objective.

\subsection{Summary and Decision}

Under the pre-declared stopping conditions (stop if progress is too slow;
otherwise cap the run at 1 hour), this branch was closed:

\begin{quote}
\textit{STOP: The direct pairwise-product MILP has a weak LP relaxation
that did not certify global optimality in the observed run.
No response-aware allocation or improvement claim follows from the
attempted runs.}
\end{quote}

This outcome does not imply that response-aware allocation with exact
quota constraints is mathematically intractable. A dedicated MIQP solver
(e.g., Gurobi or CPLEX) or a stronger relaxation strategy may
certify global optimality at 48-tap scale. This remains future work
conditioned on solver availability.
